\SourceLeaf{313}
\[
\begin{aligned}
\log b={}&\log\beta+\log m\\
&-\tfrac1{24}mm(2\beta\beta-3\tau\tau)\\
&+\tfrac1{2880}m^4(26\beta^4-60\beta\beta\tau\tau+45\tau^4)\\
&-\tfrac1{362880}m^6(502\beta^6-1638\beta^4\tau\tau+1890\beta\beta\tau^4-945\tau^6)\\
&+\tfrac1{29030400}m^8(7102\beta^8-30120\beta^6\tau\tau+49140\beta^4\tau^4\\
&\qquad\qquad-37800\beta\beta\tau^6+14175\tau^8)\\
&-\text{etc.}
\end{aligned}
\]

\indent Substituting the values given above for \(\log m,mm,m^4\) and so on gives the series sought. It can also be derived directly from the first series, in powers of \(\beta,\lambda,r\), by substituting \(rr-\beta\beta-\tau\tau\) for \(\lambda\lambda\).

\Article{19}
For our purpose, it is enough to give the formulas accurately through the fourth order inclusive, as follows.
\[
\begin{aligned}
\log t={}&\log\tau+\tfrac1{24}(2rr+\tau\tau)+\tfrac1{2880}(14r^4+20rr\tau\tau+11\tau^4)\\
\log l={}&\log\lambda-\tfrac1{24}(rr-\lambda\lambda)-\tfrac1{2880}(r^4+10rr\lambda\lambda-11\lambda^4)\\
\log b={}&\log\beta-\tfrac1{24}(rr-\beta\beta-3\lambda\lambda)\\
&-\tfrac1{2880}(r^4+10rr\beta\beta+30rr\lambda\lambda-11\beta^4-30\beta\beta\lambda\lambda-45\lambda^4)
\end{aligned}
\]
Any one of the following formulas may also be used in place of the last one:
\[
\begin{aligned}
\log b={}&\log\beta+\tfrac1{24}(2rr-2\beta\beta+3\tau\tau)\\
&+\tfrac1{2880}(14r^4-40rr\beta\beta+60rr\tau\tau+26\beta^4-60\beta\beta\tau\tau+45\tau^4)\\
\log b={}&\log\beta+\tfrac1{24}(2\lambda\lambda+\tau\tau)\\
&-\tfrac1{2880}(12\beta\beta\lambda\lambda-12\beta\beta\tau\tau-14\lambda^4-32\lambda\lambda\tau\tau+\tau^4)\\
\log b={}&\log\beta+\tfrac1{24}(2\lambda\lambda+\tau\tau)\\
&-\tfrac1{2880}(12rr\lambda\lambda-12rr\tau\tau-26\lambda^4-8\lambda\lambda\tau\tau-11\tau^4)
\end{aligned}
\]
In all these formulas, \(r,\beta,\lambda,\tau,b,l,t\) are understood to be expressed in units of the radius, and the logarithms are natural logarithms. If instead those seven quantities are to be expressed in arcseconds and Briggs's logarithms are to be used, the formulas require only two changes. The common numerical coefficient of the second-order terms, \(\frac1{24}\), must be replaced by \(\frac12\mu\), and the common numerical coefficient of the fourth-order terms, \(\frac1{2880}\), by \(\nu\), where \(\mu,\nu\) denote the products obtained by multiplying the quantities \(\frac1{12}\rho\rho\) and \(\frac1{2880}\rho^4\)
